\section{Introduction}

The cochlea converts acoustic vibration into neural signals through coordinated sensory, neuronal, vascular and supporting-cell functions. Acoustic overexposure perturbs this multicellular system and initiates transcriptional responses that extend beyond sensory hair cells. Experimental noise exposure can leave persistent cochlear synaptic and neural injury after hearing thresholds recover, with differential vulnerability among auditory nerve fibers~\citep{Kujawa2009,Furman2013}. Acoustic injury also induces cochlear cytokine expression and recruits mononuclear phagocytes, linking tissue stress with local immune responses~\citep{Fujioka2006,Hirose2005}.

Single-cell atlases resolve the molecular diversity of cochlear supporting, neural and lateral-wall populations across postnatal maturation and aging~\citep{Jean2023,Sun2023}. Spiral ganglion neurons comprise molecularly distinct subclasses with different afferent connectivity~\citep{Petitpre2018}. Single-cell profiling has established distinct neuronal and lateral-wall responses to acoustic trauma, including stress-response programs and changes in ion-transport genes~\citep{Milon2021}. Whole-cochlea proteomic and transcriptomic analyses have also connected noise exposure with proteotoxic stress and recruitment of the proteostasis network~\citep{Jongkamonwiwat2020}.

A central question is how a tissue-level inflammatory signature is distributed among the cells that contribute to it. A pathway can be preferentially expressed in one cell population while changing in different directions across populations after exposure. Earlier reporter studies identified noise-associated nuclear factor \(\kappa\)B (NF-\(\kappa\)B) activity in cochlear fibrocytes and linked the response to cell type and mouse strain~\citep{Adams2009,Herranen2018}. Replicated transcriptomic datasets now permit sample-level estimates of these cellular expression responses. Such estimates retain the biological sample as the unit of inference, following established frameworks for replicated single-cell differential-state analysis~\citep{Crowell2020,Squair2021}.

We asked whether cochlear cell populations contributing to inflammatory gene expression showed similar or opposing responses after acoustic trauma. GSE168041 provided replicated spiral ganglion and lateral-wall profiles for estimating cell-type-specific responses~\citep{Milon2021}. GSE160639 provided whole-cochlea measurements across three sound levels~\citep{Jongkamonwiwat2020}, and GSE252285 provided expression localization across control, hidden hearing loss and noise-induced hearing loss conditions~\citep{GSE252285}. Animal-level pathway estimates and leave-one-animal-out analyses resolve cell-type inflammatory divergence, complemented by annotation sensitivity of the expression-localization patterns.
